\clearpage
\fancyhead[L]{\small CARLINI. EFFEMERIDI ASTRONOMICHE.}
\hypertarget{native-571-lower}{}
\begin{center}\small Göttingische gelehrte Anzeigen. Stück 48. Seite 475..478. 1814 März 24.\\[3pt]\rule{.72\linewidth}{.3pt}\end{center}

{\rightskip=0pt plus 3em \emph{Effemeridi astronomiche di Milano per l'anno} 1814, \emph{calcolate da} \textsc{Francesco Carlini}. \emph{Dalla stamperia reale. Con appendice} 120 und 140 Seiten in klein Quart.\par}

Da wir von der zweckmässigen Einrichtung dieser vortrefflichen astronomischen Ephemeriden und von den \emph{stehenden} Artikeln, welche ihm beigegeben sind, bereits bei der Anzeige der vorhergehenden Jahrgänge Rechenschaft gegeben haben (m.\,s. St. 95 vom vorigen Jahre [S. 556 d.\,B.]), so schränken wir uns jetzt auf die Anzeige des auch bei diesem Jahrgange sehr reich ausgestatteten Anhangs ein. Wir finden zuerst die Beobachtungen des ersten Cometen des Jahrs 1811 von \textsc{Barnabas Oriani}. Diese Beobachtungen reichen vom 29. August 1811 bis 21. Januar 1812; es werden sowohl die Originalvergleichungen am Aequatorealsector als die daraus abgeleiteten Stellungen mitgetheilt, auch parabolische Elemente, wodurch diese sehr gut, aber nicht so genau die ältern \textsc{von Zach}'schen Beobachtungen dargestellt werden. Was \textsc{Oriani} über das äussere Ansehen dieses merkwürdigen Cometen sagt, ist übereinstimmend mit den Beobachtungen anderer Astronomen; er stellt auch eine Vergleichung mit dem Cometen von 1744 an, welcher in dieser Beziehung mit dem von 1811 viel ähnliches hatte, allein mit Recht bemerkt dieser einsichtsvolle Astronom, dass die Anzahl grösserer mit Fernröhren beobachteter Cometen noch viel zu klein, uns es daher zu misslich sei, auf einige bei drei oder vier Cometen im äussern Ansehen wahrgenommene Aehnlichkeit Theorien gründen zu wollen, die auf eine plausible Art den Ursprung der verschiedenen seltsamen Erscheinungen ihrer Atmosphären und Schweife erklären könnten. Am 24. December glaubte \textsc{Oriani} \emph{mitten} in dem Kopfnebel des Cometen einen hellen Kern zum ersten Male zu bemerken, allein da derselbe in den folgenden Tagen sich nicht wieder unterscheiden liess, so wurde er hierüber wieder zweifelhaft und hielt es für wahrscheinlich, dass jener vermeinte Kern vielmehr ein Fixstern 9. oder 10. Grösse gewesen sei. Es scheint uns wohl der Mühe werth, dass deshalb an dem Platze, wo der Comet den 24. December stand, (ger. Aufst. $306^\circ\,16'\,21''$, Abw. $1^\circ\,54'\,6''$ Nordl.) einmal wieder nachgesehen werde (in der gegenwärtigen Jahreszeit ist dieser Theil des Himmels nicht sichtbar). Hierauf folgen in gleichem Detail die Beobachtungen des zweiten Cometen von 1811, von demselben Astronomen, auch parabolische Elemente, welche indessen keine sehr gute Uebereinstimmung geben. Bekanntlich hat Herr \textsc{Nicolai} die Bahn dieses Cometen in einer Ellipse berechnet, wodurch auch die \textsc{Oriani}'schen Beobachtungen, ohne vorher mit benutzt gewesen zu sein, doch recht gut dargestellt werden. --- Beobachtungen des Cometen vom Jahre 1812 von \emph{Ebendemselben} (vom 1. bis 25. Sept.) nebst parabolischen Elementen. --- Es folgen hierauf in einer Reihe von Aufsätzen mehrere beobachtete Planetenpositionen von \textsc{Oriani} und \textsc{Santini}, die wir hier nur kurz berühren können. Von \textsc{Oriani} sind die Gegenscheine des Uranus 1811, des Mars 1811, der Vesta 1811 und der Vesta 1812, die drei erstern am \textsc{Ramsden}'schen Mauerquadranten, die vierte am Aequatorealsector; von \textsc{Santini} in Padua die Opposition der Juno 1810 und die des Uranus 1810, beide am \textsc{Ramsden}'schen Mauerquadranten. --- \SourceLeaf{572}Sternbedeckungen in den Jahren 1812 und 1813 und Sonnenfinsterniss vom 31. Januar 1813, beobachtet von \textsc{Oriani}, nebst einer ausführlichen Berechnung der Bedeckung von $\alpha$ Stier 22. October 1812, um die Fehler der neuen \textsc{Burckhardt}'schen Mondstafeln zu bestimmen. Es wird manchem lieb sein, hier die nöthigen Vorschriften in zierlicher Form zusammengestellt und practisch erläutert zu sehen; vielleicht wäre indess für manche Leser die Bemerkung nicht ganz überflüssig gewesen, dass eben in dem gegenwärtigen Fall das Resultat für den Breitenfehler der Tafeln nicht sehr scharf ausfallen konnte. --- Ein Brief des Hrn. \textsc{Joseph Piazzi} an \textsc{Oriani} vom 4. Juli 1812, welcher noch einige Verbesserungen seines grossen Sterncatalogs und des kleinern im Libro Sesto della specola astronomica, und ausserdem die angenehme Nachricht enthält, dass jener berühmte Beobachter eine ganz neue Bearbeitung des grossen Catalogs der Vollendung schon ganz nahe gebracht hat, die die Astronomen nicht anders als mit Ungeduld erwarten können. --- Betrachtungen über die astronomischen Uhren von \textsc{Angelo Cesaris} enthalten neben dem Bekannten einige nicht uninteressante Bemerkungen über verschiedene von Italienischen Künstlern versuchte Abänderungen des Echappements und anderer Theile, welchen Bemerkungen man nur etwas mehr Ausführlichkeit wünschen möchte. --- Noch einige beobachtete und berechnete Oppositionen der neuen Planeten von \textsc{Francesco Carlini}, nemlich Beobachtungen der Ceres 1811 und 1812 und der Pallas 1811, und Berechnung der Oppositionen der Ceres von 1809, 1811 und 1812, der Vesta 1810, der Juno 1810 und der Pallas 1811. Alle diese Resultate, so wie die oben angeführten, sind, da sie von sehr geübten Beobachtern herrühren und sich auf Beobachtungen mit vortrefflichen Werkzeugen gründen, für die Planetentheorie von sehr hohem Werthe. --- Den Schluss machen die dreijährigen von \textsc{Cesaris} auf der Mailänder Sternwarte angestellten meteorologischen Beobachtungen (1810---1812), worüber wir hier uns nur die einzige Bemerkung erlauben, dass der in den nördlichen Gegenden von Europa so ungewöhnlich heisse Sommer von 1811 es in Mailand in bei weitem geringerm Grade war; während der letzten Hälfte des Juli war der höchste Thermometerstand nur $25^\circ\,4$, und plötzliche Abänderungen hatten gar nicht statt, da hingegen z. B. in Göttingen das Thermometer im Schatten den 19. Juli auf $27^\circ$, und den 20. sogar auf $29^\circ$ stieg, und dann auf einmal sehr tief herunter ging: eine niederschlagende Bemerkung für die Liebhaber allgemeiner Wettertheorien!
